% !TEX program = xelatex
%%============================================================================
%\documentclass[handout,169,8pt]{beamer}
%\usepackage{pgfpages,tikz}
%
%\pgfpagesuselayout{6 on 1}[
%letterpaper,
%border shrink=-5mm
%]
%
%\mode<handout>{
%    \setbeamertemplate{background}{
%        \begin{tikzpicture}[overlay, remember picture]
%            \draw [line width=0.8pt, color=gray]
%            ([shift={(0.1cm,-0.1cm)}]current page.north west)
%            rectangle
%            ([shift={(-0.1cm,0.1cm)}]current page.south east);
%        \end{tikzpicture}
%    }
%}
%%============================================================================
\documentclass[8pt,aspectratio=169]{beamer}

% Custom Beamer settings (your local style file)
\RequirePackage{./.xbmr}

% Course Information
\title[EDA 898]{EDA 898}
\subtitle{Reviews of Linear Algebra and Matrix Computations}
\author{Dr.~Behzad Nouri}
\date{Winter 2026}

\begin{document}

% In handout mode, the title page is essential for the first slide slot
%%%%%%%%%%%%%%%%%%%%%%
%%%     Slide      %%%
%%%%%%%%%%%%%%%%%%%%%%
\begin{frame}
  \titlepage
\end{frame}


%%%%%%%%%%%%%%%%%%%%%%
%%%     Slide      %%%
%%%%%%%%%%%%%%%%%%%%%%
\begin{frame}[fragile,t]{Prologue}
  In the next slides, we will review some fundamental concepts in linear algebra and matrix analysis that are essential for understanding the mathematical foundations of EDA. We will cover topics such as eigenvalues, eigenvectors, diagonalization and matrix functions as they are needed in this EDA course. These concepts will provide the necessary background for analyzing and designing complex systems in the context of electronic design automation.
\end{frame}


%%%%%%%%%%%%%%%%%%%%%%%%%%%%%%%%%%%%%%%%%%%%%%%%%%%%%%%%%
%%%   Sections:
%%%%%%%%%%%%%%%%%%%%%%%%%%%%%%%%%%%%%%%%%%%%%%%%%%%%%%%%%%
%%%%%%%%%%%%%%%%%%%%%%%%%%%%%%%%%%%%%%%%%%%%%%%%%
%%  Section: Eigenvalues, Eigenvectors, and Diagonalization
%%%%%%%%%%%%%%%%%%%%%%%%%%%%%%%%%%%%%%%%%%%%%%%%%
\section{Eigenvalues, Eigenvectors, and Diagonalization}

%%%%%%%%%%%%%%%%%%%%%
%%     Slide      %%%
%%%%%%%%%%%%%%%%%%%%%    
%%Note: 'fragile' is mandatory when using listings in a frame
\begin{frame}[fragile,t,allowframebreaks]{Eigenvalues and Eigenvectors --}

    \begin{alertblock}{\textbf{\underline{Eigenvectors and Eigenvalues:}}} \medskip
        An \alert{eigenvector} of an $n \times n$ matrix $\mathbf{A}$ is a nonzero vector $\mathbf{v}$ such that
        $$\mathbf{A}\mathbf{v} = \lambda \mathbf{v}$$
        for some scalar $\lambda$.

        \vspace{0.3cm}

        $\lambda$ is an \alert{eigenvalue} of $\mathbf{A}$ if there exists a nontrivial solution $\mathbf{v}$ ($\mathbf{v} \neq \mathbf{0}$) to this equation; such a vector $\mathbf{v}$ is called an \alert{eigenvector corresponding to $\lambda$}.
    \end{alertblock}

    From the above definition, a scalar $\lambda$ is an eigenvalue of an $n \times n$ matrix $\mat A$ iff the homogeneous linear system
    %
    \begin{equation}
        \left(\mat{A} - \lambda \mat{I}\right)\,\mathbf{v} = \mathbf{0}
    \end{equation}
    %
    has a nontrivial solution ($\mathbf{v} \neq \mathbf{0}$). Since $\mathbf{v}$ is nonzero, the matrix $\left(\mat{A} - \lambda \mat{I}\right)$ must be singular, i.e., \emph{non-invertible}. A matrix is singular if and only if its determinant is zero; therefore, the eigenvalues of $\mat A$ are precisely the scalars $\lambda$ for which

    \begin{equation}
        \det\left(\mat{A} - \lambda \mat{I}\right) = 0\,.
    \end{equation}

    \begin{alertblock}{\textbf{\underline{Characteristic Polynomial:}}}\medskip
        For a square matrix $\mathbf{A}$, $\det\left(\mat{A} - \lambda \mat{I}\right)$ is a polynomial of degree $n$ called the \alert{characteristic polynomial} of $\mat A$, i.e.,
        \[
            p(\lambda)
            = \det(\lambda\mathbf{I} - \mathbf{A})
            = \lambda^n + b_{n-1}\lambda^{n-1} + \cdots + b_1\lambda + b_0.
        \]
    \end{alertblock}


    \begin{alertblock}{\textbf{\underline{Computing Eigenvalues:}}}\medskip
        The eigenvalues of a square matrix $\mathbf{A}$, are obtained by computing the roots of $\det\left(\mat{A} - \lambda \mat{I}\right)$. Since a degree-$n$ polynomial has exactly $n$ roots (counting multiplicities), an $n \times n$ matrix can have $n$ eigenvalues.
        \[
            p(\lambda)
            = \det(\lambda\mathbf{I} - \mathbf{A})
            = \lambda^n + b_{n-1}\lambda^{n-1} + \cdots + b_1\lambda + b_0.
        \]
    \end{alertblock}

    \bigskip
    \begin{alertblock}{\textbf{\underline{In a more mathematical setting}}}\medskip
        One can show that the set of eigenvalues of $\mat A$ is precisely the \textbf{spectrum} of $\mat A$, denoted $\sigma(\mat A)$.
        An eigenvector is a dimension that remains invariant under pre-multiplication by $\mat A$.
        \medskip
        Let $\mathcal{S}$ be a subspace as $\mathcal{S} \subseteq \mathbb{C}^n$ with the property
        \[
            \mat{v} \in \mathcal{S} \;\;\Longrightarrow\;\; \mat{A v} \in \mathcal{S}
        \]
        Hence $\mathcal{S}$ is called an \alert{invariant subspace} of $\mat A$. The importance of eigenvector is that it is spanning this invariant subspace of $\mat A$.
    \end{alertblock}

    \bigskip
    \begin{itemize}
        \item If $\lambda \in \mathbb{R},\ \lambda > 0$:
              $\mathbf{v}$ and $\mat{A}\mathbf{v}$ point in the same direction.
        \item If $\lambda \in \mathbb{R},\ \lambda < 0$:
              $\mathbf{v}$ and $\mat{A}\mathbf{v}$ point in opposite directions.
        \item If $\lambda \in \mathbb{R},\ |\lambda| < 1$:
              $\mat{A}\mathbf{v}$ is shorter than $\mathbf{v}$.
        \item If $\lambda \in \mathbb{R},\ |\lambda| > 1$:
              $\mat{A}\mathbf{v}$ is longer than $\mathbf{v}$.
    \end{itemize}

\end{frame}



%%%%%%%%%%%%%%%%%%%%%%
%%%     Slide      %%%
%%%%%%%%%%%%%%%%%%%%%%
\begin{frame}[fragile,t]{Question}
    \textbf{\red{What does it mean for a matrix $\mat{A}$ to have an eigenvalue equal to $0$?}}

    \begin{alertblock}{\textbf{\underline{Answer:}}}\medskip
        A matrix $\mat{A}$ has $0$ as an eigenvalue exactly when the equation
        \begin{equation} \label{eq:eig-3}
            \mat{A}\,\mathbf{v} = 0\,\mathbf{v}
        \end{equation}
        admits a nontrivial solution.
        Since \eqref{eq:eig-3} reduces to
        \[
            \mat{A}\,\mathbf{v} = \mathbf{0},
        \]
        a nontrivial solution exists precisely when $\mat{A}$ is \emph{not invertible}.\\[1mm]

        \begin{center}
            \textbf{$\mat{A}$ has $0$ as an eigenvalue \quad$\Longleftrightarrow$\quad $\mat{A}$ is singular.}
        \end{center}
    \end{alertblock}

\end{frame}


%%%%%%%%%%%%%%%%%%%%%
%%     Slide      %%%
%%%%%%%%%%%%%%%%%%%%% 
% Note: 'fragile' is mandatory when using listings in a frame
\begin{frame}[t,fragile,allowframebreaks]{Example: Eigen Decomposition of a Matrix --}
    To compute eigenvalues:
    \begin{equation*}
        \mat{A}\mat{v}\:=\:\lambda\mat{v} \quad \Longrightarrow \left(A-\lambda\mat{I}\right)\mat{v} = \mat{0}
    \end{equation*}
    To have a nontrivial solution as $\mat{v}\neq\mat{0}$, it is,
    \begin{equation*}
        \det\left(A-\lambda\mat{I}\right) = \left|A-\lambda\mat{I}\right| =0
    \end{equation*}
    \medskip
    \textbf{Example:}

    \[
        A =
        \begin{bmatrix}
            3 & 2 \\
            1 & 2
        \end{bmatrix}
        \qquad
        |A - \lambda I| = 0
        \;\;\Longrightarrow\;\;
        \lambda^{2} - 5\lambda + 4 = 0
        \;\;\Longrightarrow\;\;
        \text{Eigenvalues: } 1,\; 4
    \]

    One can also investigate that: $\mat{A}^{2} - 5\mat{A} + 4\mat{I} = \mat{0}.$
    \vfill
    \begin{equation*}
        \mat{A}\mat{v}_1=\lambda_1\mat{v}_1\;\;\Longrightarrow\;\;
        \begin{bmatrix}
            3 & 2 \\
            1 & 2
        \end{bmatrix} \begin{bmatrix}x_1\\x_2\end{bmatrix}  = 1\times
        \begin{bmatrix}v_1\\varbigtriangledown_2\end{bmatrix}
    \end{equation*}

    \begin{equation*}
        \begin{cases}
            2v_1+2v_2=0 \\
            v_1+v_2=0
        \end{cases}
        \;\; \Longrightarrow\;\; v_1+v_2=0\;\;\text{for }v_1=1 \;\; \Longrightarrow\;\;v_2=-1
    \end{equation*}
    \begin{center}
        One eigen-pair is   ($\lambda_1=1$, $\mat{x}_1=\begin{bmatrix}1\\-1\end{bmatrix}$)
    \end{center}
    and
    \begin{equation*}
        \mat{A}\mat{v}_2=\lambda_2\mat{v}_2\;\;\Longrightarrow\;\;
        \begin{bmatrix}
            3 & 2 \\
            1 & 2
        \end{bmatrix}
        \begin{bmatrix}v_1\\v_2\end{bmatrix}
        = 4\times   \begin{bmatrix}v_1\\v_2\end{bmatrix}
    \end{equation*}

    \begin{equation*}
        \begin{cases}
            -v_1+2v_2=0 \\
            v_1-2v_2=0
        \end{cases}
        \;\; \Longrightarrow\;\; -v_1+2v_2=0\;\;\text{for }v_1=2 \;\; \Longrightarrow\;\;v_2=1
    \end{equation*}

    \begin{center} Second eigen-pair is   ($\lambda_2=4$, $\mat{v}_2=\begin{bmatrix}2 \\1
            \end{bmatrix}$)
    \end{center}

    It is
    \begin{equation*}
        \begin{bmatrix}
            3 & 2 \\
            1 & 2
        \end{bmatrix}
        \begin{bmatrix}
            1  & 2 \\
            -1 & 1
        \end{bmatrix}\:=\:
        \begin{bmatrix}
            1  & 2 \\
            -1 & 1
        \end{bmatrix}
        \begin{bmatrix}
            1 & 0 \\
            0 & 4
        \end{bmatrix}
    \end{equation*}

    In a general notation it is
    \begin{equation*}
        \mat{A}\mat{V}\:=\:\mat{V}\mat{\Lambda}
    \end{equation*}
    \medskip
    where $\mat{V}$ is the matrix of eigenvectors and $\mat{\Lambda}$ is the diagonal matrix of eigenvalues.


    $\mat{T}$ being invertible, we can obtain

    \begin{equation*}
        \begin{bmatrix}
            1  & 2 \\
            -1 & 1
        \end{bmatrix}^{-1}\:
        \begin{bmatrix}
            3 & 2 \\
            1 & 2
        \end{bmatrix}
        \begin{bmatrix}
            1  & 2 \\
            -1 & 1
        \end{bmatrix}\:=\:
        \begin{bmatrix}
            1 & 0 \\
            0 & 4
        \end{bmatrix}
    \end{equation*}
    This is the process of diagonalization, where we have transformed the matrix $\mat{A}$ into a diagonal matrix $\mat{\Lambda}$ using the similarity transformation defined by $\mat{T}$.
    \begin{equation*}
        \mat{T}^{-1}\mat{A}\mat{T}\:=\:\mat{\Lambda}
    \end{equation*}
    \vfill


    \pagebreak

    \begin{alertblock}{\textbf{\underline{Interpretation:}}}\medskip
        It express that the action of $\mat{A}$ on the eigenvectors $\mathbf{v}_i$
        \[\mat{A}\mathbf{v}_i = \lambda_i\mathbf{v}_i\]
        is simply to scale them by their corresponding eigenvalues $\lambda_i$. In other words, each eigenvector $\mathbf{v}_i$ is an invariant direction under the transformation defined by $A$, and the eigenvalue $\lambda_i$ indicates how much $A$ stretches or compresses along that direction.

        \medskip
        The matrix equation can be written as:
        \[
            A\,\left[\, \mathbf{v}_1 \;\; \cdots \;\; \mathbf{v}_n \,\right]
            =
            \left[\, \mathbf{v}_1 \;\; \cdots \;\; \mathbf{v}_n \,\right]
            \begin{bmatrix}
                \lambda_1             \\
                 & \ddots             \\
                 &        & \lambda_n
            \end{bmatrix}.
        \]
    \end{alertblock}

\end{frame}

%%%%%%%%%%%%%%%%%%%%%%
%%%     Slide      %%%
%%%%%%%%%%%%%%%%%%%%%%
\begin{frame}[t, fragile]{Distinct eigenvalues}

    \begin{alertblock}{\textbf{\underline{Fact:}}}\medskip
        If $A$ has distinct eigenvalues, then $A$ is diagonalizable.
    \end{alertblock}

    \begin{itemize}
        \item \alert{Distinct eigenvalues} means $\lambda_i \neq \lambda_j$ for $i \neq j$.

        \item The converse is \alert{not} true:
              a matrix can have repeated eigenvalues and still be diagonalizable.
    \end{itemize}

\end{frame}


%%%%%%%%%%%%%%%%%%%%%
%%     Slide      %%%
%%%%%%%%%%%%%%%%%%%%%
% Note: 'fragile' is mandatory when using listings in a frame
\begin{frame}[t,fragile]{Diagonalization}
    Given an $n \times n$ square matrix $\mathbf{A}$ with $n$ distinct eigenvalues.

    \[
        \mathbf{A}\,\mathbf{T} \;=\; \mathbf{T}\,\mathbf{\Lambda}
    \]

    If $\mat{T}$ is full rank (i.e., invertible), then $A$ is diagonalizable, We can obtain:
    \[
        \mathbf{T}^{-1}\,\mathbf{A}\mathbf{T} \:=\: \mathbf{\Lambda}\:=\:\operatorname{diad}\left(\lambda_1,\lambda_2,\dots,\lambda_n\right)
    \]
    Hence the name \textbf{Diagonalization}.

\end{frame}

%%%%%%%%%%%%%%%%%%%%%%
%%%     Slide      %%%
%%%%%%%%%%%%%%%%%%%%%%
\begin{frame}{Not all matrices are diagonalizable}

    \begin{alertblock}{Example:}\medskip
        Given the matrix  \hspace*{9pc}
        $ A =
            \begin{bmatrix}
                0 & 1 \\
                0 & 0
            \end{bmatrix}
        $
    \end{alertblock}

    \begin{itemize}
        \item characteristic polynomial:\hspace*{12pt}
              $P(\lambda) = \lambda^{2} \quad \Rightarrow \quad \lambda = 0 \text{ is the only eigenvalue}$

        \item eigenvectors satisfy \(A\mathbf{v} = 0\mathbf{v} = \mathbf{0}\), i.e.,
              \[
                  \begin{bmatrix}
                      0 & 1 \\
                      0 & 0
                  \end{bmatrix}
                  \begin{bmatrix}
                      v_1 \\
                      v_2
                  \end{bmatrix}
                  =
                  \begin{bmatrix}
                      0 \\
                      0
                  \end{bmatrix}
              \]

        \item so all eigenvectors must satisfy \(v_2 = 0\), hence
              \[
                  \mathbf{v} =
                  \begin{bmatrix}
                      v_1 \\[2pt]
                      0
                  \end{bmatrix},
                  \qquad v_1 \neq 0
              \]

        \item thus, $\mat A$ has only one linearly independent eigenvector

        \item therefore, $\mat A$ is \textbf{not diagonalizable}
        \item if $\mat A$ is not diagonalizable, it is sometimes called \textbf{\red{defective}}
    \end{itemize}
    \vfill
\end{frame}




%%%%%%%%%%%%%%%%%%%%%%
%%%     Slide      %%%
%%%%%%%%%%%%%%%%%%%%%%
\begin{frame}[t, fragile]{Diagonalization of Symmetric Real Matrices}
    Let $\mathbf{A}_{n\times n}$ be a real symmetric matrix, i.e., $\mathbf{A}^T = \mathbf{A}$.
    Then $\mathbf{A}$ admits an orthogonal diagonalization:
    \[
        \mathbf{A} = \mathbf{V}\,\mathbf{\Lambda}\,\mathbf{V}^{T},
    \]
    where $\mathbf{T}^{T} = \mathbf{T}^{-1}$ and
    \[
        \mathbf{T}^{T}\mathbf{T} = \mathbf{T}\mathbf{T}^{T} = \mathbf{I}.
    \]
    Here, $\mathbf{T}$ contains the orthonormal eigenvectors of $\mathbf{A}$, and
    $\mathbf{\Lambda}$ is the diagonal matrix of eigenvalues.

\end{frame}


%%%%%%%%%%%%%%%%%%%%%%
%%%     Slide      %%%
%%%%%%%%%%%%%%%%%%%%%%   
\begin{frame}<0>{Invariant Subspaces and Similarity}

    \textbf{The Fundamental Relation}
    \vspace{0.2cm}
    Let $A \in \mathbb{C}^{n \times n}$ and $B \in \mathbb{C}^{k \times k}$. Consider the matrix equation:
    \[ AX = XB, \quad X \in \mathbb{C}^{n \times k} \]

    \hrule
    \vspace{0.4cm}

    \begin{block}{1. Invariant Subspaces}
        The range of $X$, denoted $\text{ran}(X)$, is an \textbf{invariant subspace} under $A$.
        If $v \in \text{ran}(X)$, then $Av \in \text{ran}(X)$.
    \end{block}

    \begin{block}{2. Eigenvalue Relationship}
        If $(\lambda, y)$ is an eigenpair of $B$ (i.e., $By = \lambda y$), then:
        \[ A(Xy) = X(By) = \lambda(Xy) \]
        If $Xy \neq 0$, then $\lambda$ is an eigenvalue of $A$ with eigenvector $Xy$.
    \end{block}

\end{frame}


%%%%%%%%%%%%%%%%%%%%%%
%%%     Slide      %%%
%%%%%%%%%%%%%%%%%%%%%%  
\begin{frame}[t,fragile]{Similarity Transformations}
    %   
    \begin{itemize}
        \item \textbf{Similarity:}
              If $Q$ is \textbf{square} and \textbf{nonsingular}, then
              \[
                  B = Q^{-1} A Q .
              \]
              In this case, $A$ and $B$ are said to be \textbf{similar}
              (denoted $A \sim B$).

              \medskip

        \item \textbf{Spectrum Inclusion:}
              If $Q$ has \textbf{full column rank} (but is not necessarily square),
              then every eigenvalue of $B$ is an eigenvalue of $A$:
              \[
                  \lambda(B) \subseteq \lambda(A).
              \]

              \medskip

        \item \textbf{Spectral Equality:}
              If $Q$ is \textbf{invertible} (square and nonsingular), then similar matrices
              share exactly the same eigenvalues:
              \[
                  \lambda(A) = \lambda(B).
              \]

    \end{itemize}

    \begin{alertblock}{Takeaway:}\smallskip
        The matrix $Q$ acts as a \textbf{change of basis}, translating the
        representation of $B$ into the coordinate system of $A$ while preserving
        its spectral structure.
    \end{alertblock}
    \vfill

\end{frame}


%%EoF

%%%%%%%%%%%%%%%%%%%%%%%%%%%%%%%%%%%%%%%%%%%%%%%%%
%%  Section: Matrix Functions
%%%%%%%%%%%%%%%%%%%%%%%%%%%%%%%%%%%%%%%%%%%%%%%%%
\section{Matrix Functions}

%%%%%%%%%%%%%%%%%%%%%%
%%%     Slide      %%%
%%%%%%%%%%%%%%%%%%%%%% 
\begin{frame}[t]{Matrix Functions}
    For a diagonalizable matrix $A = T \Lambda T^{-1}$, we can extend a scalar function $f$ (such as polynomials, exponentials, or roots) to the matrix case.

    \begin{alert}
        If $f$ is defined on the spectrum%%%%%%%%%%%%
        \footnote{The \emph{spectrum} of a matrix $A$, denoted $\sigma(A)$, is the set of all eigenvalues of $A$, i.e.,
            \[
                \sigma(A) = \{\lambda \in \mathbb{C} : \det(A - \lambda I) = 0\}.
            \]}%
        of $A$, we define the \textbf{matrix function} as:
        \[
            f(A) = T \, f(\Lambda) \, T^{-1}
        \]
    \end{alert}


    where $f(\Lambda)$ is the diagonal matrix obtained by applying $f$ to each eigenvalue:
    \[
        f(\Lambda) = \operatorname{diag}\left(f(\lambda_1), f(\lambda_2), \dots, f(\lambda_n)\right)
    \]


    \vfill
\end{frame}


%%%%%%%%%%%%%%%%%%%%%%
%%%     Slide      %%%
%%%%%%%%%%%%%%%%%%%%%% 
\begin{frame}[t,allowframebreaks]{Matrix Powers --}

    \textbf{Why diagonalize?}
    Computing powers $A^n$ directly requires $n$ matrix multiplications.
    But if $A$ is diagonalizable, we can write
    \[
        A = T \Lambda T^{-1},
    \]
    and therefore
    \[
        A^n = (T \Lambda T^{-1})^n = T\, \Lambda^n\, T^{-1}.
    \]

    Since $\Lambda$ is diagonal, its powers are trivial:
    \[
        \Lambda^n =
        \mathrm{diag}(\lambda_1^n, \lambda_2^n, \dots, \lambda_n^n),
    \]
    so the computation becomes extremely inexpensive.

    \medskip
    \pagebreak

    \textbf{Examples:}
    \[
        A^{1/2} = T\, \Lambda^{1/2}\, T^{-1},
        \qquad
        e^A = T\, e^{\Lambda}\, T^{-1},
    \]
    with
    \[
        e^\Lambda =
        \mathrm{diag}(e^{\lambda_1},\, \dots,\, e^{\lambda_n}).
    \]

    \begin{alertblock}{Why this matters (especially in EDA):}\smallskip
        Diagonalization converts difficult matrix computations into simple scalar
        operations on eigenvalues.
        This enables fast evaluation of $A^n$, $A^{1/2}$, and $e^A$, which arise in
        simulation algorithms, stability analysis, and time‑domain state‑transition solutions.
    \end{alertblock}

    \vfill

\end{frame}

\begin{frame}[t, allowframebreaks]{Matrix Exponential --}

    \textbf{Scalar exponential (review):}
    \[
        e^{x} = 1 + x + \frac{x^{2}}{2!} + \frac{x^{3}}{3!} + \cdots
    \]

    \medskip

    \textbf{Matrix exponential:}
    For any square matrix $A$, the exponential is defined by the same power series:
    \[
        e^{A} \;=\; \sum_{n=0}^{\infty} \frac{A^{n}}{n!}.
    \]

    \medskip

    \textbf{Key idea:}
    The series converges for every matrix $A$, and $e^{A}$ behaves like a
    generalization of the scalar exponential to matrices.

    \medskip


    \textbf{1. Power–series definition:}
    \[
        e^{A}
        = \sum_{n=0}^{\infty} \frac{A^{n}}{n!}.
    \]

    \medskip

    \textbf{2. Diagonalization:}
    If $A$ is diagonalizable,
    \[
        A = T \Lambda T^{-1},
    \]
    where $\Lambda = \mathrm{diag}(\lambda_1,\dots,\lambda_n)$.

    \medskip

    \textbf{3. Put the series into the diagonal form:}
    \[
        A^{n} = (T \Lambda T^{-1})^{n}
        = T\,\Lambda^{n}\,T^{-1}.
    \]

    \medskip

    \textbf{4. Substitute into the series:}
    \[
        e^{A}
        = \sum_{n=0}^{\infty} \frac{A^{n}}{n!}
        = \sum_{n=0}^{\infty}
        \frac{T \Lambda^{n} T^{-1}}{n!}
        = T\left(\sum_{n=0}^{\infty} \frac{\Lambda^{n}}{n!}\right)T^{-1}.
    \]

    \medskip

    \textbf{5. But $\Lambda$ is diagonal, so}
    \[
        e^{\Lambda}
        = \mathrm{diag}(e^{\lambda_1},\, \dots,\, e^{\lambda_n}),
    \]
    and therefore
    \[
        \boxed{\,e^{A} = T e^{\Lambda} T^{-1}\, }.
    \]

\end{frame}


%%EoF



%%%%%%%%%%%%%%%%%%%%%%%%%%%%%%%%%%
%%%     Bibliography Slide     %%%
%%%%%%%%%%%%%%%%%%%%%%%%%%%%%%%%%%
\begin{frame}[allowframebreaks]{References}
  \nocite{*}
  \printbibliography
\end{frame}

\end{document}